\documentclass[11pt,a4paper]{article}
\usepackage[utf8]{inputenc}
\usepackage[T1]{fontenc}
\usepackage{amsmath, amssymb, amsfonts}
\usepackage{graphicx}
\usepackage{geometry}
\geometry{a4paper, margin=1in}
\usepackage{hyperref}
\usepackage{xcolor}

\title{\textbf{Understanding ACT: CVAE and Policy Forward Pass}}
\author{Explaining the Style Variable Generation}
\date{\vspace{-5ex}}

\begin{document}

\maketitle

\section{Introduction}
During the training phase of Action Chunking with Transformers (ACT), the model behaves as a Conditional Variational Autoencoder (CVAE). This means before the main Transformer (the Decoder) can process the vision and proprioception context, we must compute a "style variable," denoted as $z$. This variable is designed to capture the stochasticity, multi-modality, and specific human nuances of a given demonstration. 

This document details the pipeline of the CVAE Encoder up to the point where all arguments are prepared for the main loss function.

\section{The CVAE Encoder Inputs}
To calculate $z$ during training, the CVAE Encoder is fed a combination of the current state and the actual future actions from the dataset. 

In paper terminology, the full observation $o_t$ includes both images and state. The encoder uses $\bar{o}_t$, which is defined as the \textbf{observation without the images}. In practice, this means the encoder only sees the low-dimensional physical data. To save on heavy computation (avoiding ResNet passes for the encoder), the images are discarded here.

\begin{itemize}
    \item \textbf{Proprioception ($\bar{o}_{t}$):} The normalized joint positions of the single-arm state. This is the only part of the observation used by the encoder.
    \item \textbf{Future Actions ($a_{t:t+k}$):} The sequence of normalized $k$ future target joint positions that the human expert actually executed.
    \item \textbf{[CLS] Token:} As described in Section IV.C (page 6) of the original paper, a learnable "Class" token is prepended to the input sequence. Because Transformers output a sequence of the same length as their input, we need a mechanism to condense a list of $k$ actions into a \textit{single} vector. The \texttt{[CLS]} token acts as a \textbf{Designated Style Extractor}---a dummy variable whose sole purpose is to "listen" to all the other tokens and serve as the final summary bucket.
\end{itemize}

\section{The Encoder Forward Pass}
The inputs are mapped to the hidden dimension size (e.g., 512) and passed through a Transformer Encoder. The paper refers to this as "BERT-like." This simply means it is a bidirectional encoder: instead of reading actions sequentially from left to right (like predicting the next word in a sentence), it looks at the entire sequence of future actions all at once to fully understand the context.

This architecture acts as a \textbf{Universal Translator}: it reads a "paragraph" of complex human physical movements and translates it into a single "word" (the style vector $z$). This forces the model to take everything that happened in a long, complicated human demonstration and summarize it into one single spot (the final \texttt{[CLS]} token). 

\subsection{Aggregating Information}
Inside this Encoder, multi-head self-attention allow the \texttt{[CLS]} token to interact with the entire mathematical sequence of future actions and the proprioception vector. By the time the data passes through all the layers:
\begin{itemize}
    \item The output vectors corresponding to the individual actions are discarded.
    \item The single vector located at the \texttt{[CLS]} position is extracted. Because of self-attention, this single vector now contains a summarized mathematical representation of the entire human action sequence's "style" in that specific state.
\end{itemize}

\subsection{Predicting the Distribution: The Gaussian Assumption}
Instead of generating a static tensor, ACT assumes the "style" of a trajectory is a random process following a \textbf{Multivariate Gaussian Distribution}. To solve this, the summary \texttt{[CLS]} feature is passed through two linear layers to predict the statistical parameters of that distribution:
\begin{enumerate}
    \item \textbf{Mean Vector ($\mu$):} The center of the "style" distribution in the 512-dimensional space.
    \item \textbf{Variance ($\sigma^2$):} A vector representing the diagonal variance (the "uncertainty" or "spread") of each style dimension.
\end{enumerate}

During training, the model performs \textbf{Dual Learning}: it simultaneously refines its ability to "translate" human motion (the Encoder's job, using weights $\phi$) while also sharpening its ability to "execute" those movements (the Policy Transformer's job, using weights $\theta$). 

\section{Sampling and the Reparameterization Trick}
Instead of directly passing the mean as $z$, we sample from the predicted Gaussian distribution $\mathcal{N}(\mu, \sigma^2)$. 

To maintain differentiability for backpropagation during training, a concept called the \textit{reparameterization trick} is applied:
\[ z = \mu + \sigma \cdot \epsilon \]
Where $\epsilon$ is random noise sampled from a standard normal distribution $\mathcal{N}(0, 1)$. This $z$ vector is finally projected to the standard hidden dimension (e.g., 512).

\section{The Policy Transformer: Action Generation}
Now that $z$ has been sampled, the model moves to its second major component: the Policy Transformer. This is a standard Transformer Encoder-Decoder architecture responsible for physically commanding the robot.

\subsection{The Policy Encoder: Building Context}
The Policy Encoder takes the normalized sensor data and combines it with the style variable to create a unified understanding of the current "scene."
The input sequence is constructed by concatenating:
\begin{enumerate}
    \item \textbf{Vision Tokens:} The $600-900$ visual features extracted from the 2-3 cameras (plus 2D positional embeddings).
    \item \textbf{Proprioception Token:} The single token representing the robot's state.
    \item \textbf{Style Token:} The projected $z$ vector.
\end{enumerate}

Inside the Policy Encoder, self-attention layers allow all these modalities to interact. The wrist camera features can "contextualize" themselves using the joint angles, while keeping the human's style ($z$) in mind. The output is a highly enriched set of memory tokens.

\subsection{The Policy Decoder: Predicting the Chunk}
The Policy Decoder's job is to read the Encoder's memory and output a sequence of $k$ future actions (e.g., $k=100$).
\begin{itemize}
    \item \textbf{Queries:} Instead of feeding inputs sequentially like a language model, the Decoder is given $k$ fixed positional embeddings. Think of these as $100$ empty "slots" asking the Encoder, "What should action 1 be?", "What should action 2 be?", etc.
    \item \textbf{Cross-Attention:} Through cross-attention, these slots query the rich memory from the Policy Encoder, filling themselves with geometric and semantic information.
    \item \textbf{Output ($\hat{a}_{t:t+k}$):} The output of the Decoder is passed through a final linear layer to predict the $k$ normalized, continuous joint actions for the robot.
\end{itemize}

\section{Preparing for the Loss Function}
By running both the CVAE Encoder ($\phi$) and the Policy Transformer ($\theta$), we have successfully generated every mathematical argument needed to compute the two halves of the total loss:

\begin{itemize}
    \item \textbf{KL Divergence Arguments:} We have the predicted parameters ($\mu$, $\sigma^2$) from the CVAE Encoder. These are used to compute the KL loss, ensuring the style distribution doesn't stray too far from a standard normal prior $\mathcal{N}(0, I)$.
    \item \textbf{Reconstruction Arguments:} We now have the Policy Decoder's final predictions ($\hat{a}_{t:t+k}$). We compare these directly to the ground truth human actions ($a_{t:t+k}$) using an L1 smooth reconstruction loss.
\end{itemize}

\textbf{Simple Summary:} To recap the forward pass, we take the expert actions and joint state to "translate" them into a style distribution ($\mu, \sigma^2$). We sample a style vector $z$ from this distribution and feed it into the main Policy Transformer alongside the camera images. This produces a predicted chunk of 100 actions. We now have everything we need to calculate how far off our predictions were (L1 Loss) and how messy our translation was (KL Loss).

\end{document}
